\section{Problem and Complementarity Structure}\label{sec:problem}

We consider the strictly convex non-negative quadratic program~\eqref{eq:bqp},
\[
  \min_{x \geq 0}\; f(x) = \tfrac{1}{2}\,x^\top A x - b^\top x,
  \qquad A = A^\top \succ 0,\quad b \in \mathbb{R}^n,
\]
together with its two structural specialisations. In the \emph{least-squares} form
$A = M^\top M$ and $b = M^\top d$ for $M \in \mathbb{R}^{m\times n}$ of full column rank, so
that $f(x) = \tfrac{1}{2}\|Mx - d\|_2^2 - \tfrac{1}{2}\|d\|_2^2$ and~\eqref{eq:bqp} is
NNLS~\cite{lawson1974}. In the \emph{equality-augmented} form a single normalisation
$\mathbf{1}^\top x = \beta$ is imposed as well,
\begin{equation}\label{eq:eqp}
  \min_{x}\; \tfrac{1}{2}\,x^\top A x - b^\top x
  \quad\text{subject to}\quad \mathbf{1}^\top x = \beta,\; x \geq 0,
\end{equation}
which constrains $x$ to the affine slice $\Delta_\beta$ of the non-negative orthant.

\paragraph{Optimality and complementarity.}
Because $f$ is strictly convex and the feasible set of~\eqref{eq:bqp} is a non-empty closed
convex cone, a unique minimiser exists and the Karush-Kuhn-Tucker (KKT) conditions are
necessary and sufficient. Introducing a multiplier $s \geq 0$ for $x \geq 0$, stationarity of
the Lagrangian $\tfrac{1}{2}x^\top A x - b^\top x - s^\top x$ gives the reduced gradient
$s = \nabla f(x) = A x - b$, and the KKT system is the linear complementarity problem
$\mathrm{LCP}(A,\,{-b})$:
\begin{equation}\label{eq:lcp}
  s = A x - b, \qquad x \geq 0, \qquad s \geq 0, \qquad x_i\,s_i = 0 \ \ \forall i.
\end{equation}
Complementary slackness partitions the indices into a \emph{free} set
$\mathcal{F} = \{i : x_i > 0\}$, on which $s_i = 0$ and hence $(Ax)_i = b_i$, and a
\emph{bound} set $\mathcal{B} = \{i : x_i = 0\}$, on which $s_i = (Ax)_i - b_i \geq 0$. A
bound variable with $s_i < 0$ would lower $f$ if released to a small positive value; a free
variable with $x_i < 0$ is infeasible and must be pushed to the bound. These two violations
drive the primal and dual steps of Section~\ref{sec:nonneg}.

For the equality-augmented problem~\eqref{eq:eqp} the Lagrangian carries an additional term
$-\lambda(\mathbf{1}^\top x - \beta)$, so the reduced gradient becomes
$s = A x - b - \lambda\mathbf{1}$ and the free-set stationarity reads
$(Ax)_i = b_i + \lambda$ for $i \in \mathcal{F}$: every free coordinate's gradient equals the
common shadow price $\lambda$ of the normalisation. Section~\ref{sec:reduction} shows that
$\lambda$ need not be carried as an unknown in the linear solve; it is recovered in closed form
from two SPD solves.

\paragraph{The $P$-matrix property.}
Since $A \succ 0$, every principal submatrix $A_{\mathcal{F}} = A_{\mathcal{F},\mathcal{F}}$
is itself SPD and hence invertible with $\det A_{\mathcal{F}} > 0$; a matrix all of whose
principal minors are positive is a $P$-matrix~\cite{murty1988}. Two consequences follow, both
used later. The linear system $A_{\mathcal{F}}\,x_{\mathcal{F}} = b_{\mathcal{F}}$ solved on
any candidate free set is well posed, so CG applies at every outer step
(Section~\ref{sec:reduction}). And $\mathrm{LCP}(A,\,{-b})$ with a $P$-matrix has a unique
solution reachable by principal pivoting~\cite{murty1988,judice1994}, which is the backbone of
the finite-termination guarantee in Section~\ref{sec:nonneg}. When $A$ is only positive
semidefinite -- as in a rank-deficient least-squares problem, $m < n$ -- the regularisation of
Section~\ref{sec:conditioning} restores $A \succ 0$ and with it the $P$-matrix property.
